% GENERATED FILE. Edit canonical lessons, then run npm run generate:books.
% canonical-source-hash: fnv1a64:516b6345cde629dd
% canonical-lessons: ES-C555-vejez, ES-C555-anciano, ES-C555-muerte, ES-C555-muerto, ES-C555-difunto

\chapter{La vejez}
\label{ch:es-la-vejez}
\hlchaptermodality{555}

\begin{chapteropening}
\textbf{By the end of this chapter:} \emph{I can talk about old age and death: the elderly, the dead and the deceased.}

Complete the last lesson of chapter 555: I can talk about old age and death: the elderly, the dead and the deceased.
\end{chapteropening}

\section[la vejez]{la vejez — old age}

\label{lesson:ES-C555-vejez}



\begin{quote}
Before the new one: say the Spanish for the age of majority, then the Spanish for to grow old.
\end{quote}

\subsection*{You'll want to know: la vejez}
\textbf{la vejez} — \textquotedblleft{}old age\textquotedblright{}.

\begin{quote}
Vivió una vejez tranquila en el campo.
\end{quote}

\begin{quote}
\emph{He lived a quiet old age in the country.}
\end{quote}

\begin{cousinweb}
\emph{vejez} is built on \emph{viejo} 'old' with the ending \emph{-ez}. \emph{Viejo} is inherited from Latin \emph{vetulus}, from \emph{vetus} 'old' (Lewis \& Short).
\end{cousinweb}

\begin{grammarlens}[title={What you've built}]
One more word for telling, explaining and describing at B1.
\end{grammarlens}

\subsection*{Your turn}
\begin{itemize}
\raggedright
  \item \emph{Say it:} \emph{la vejez}
  \item \emph{Say it:} \emph{la vejez}, once more
  \item \emph{Say it:} a sentence of your own with \emph{la vejez}
  \item \emph{Recall:} say the Spanish for the age of majority, then the Spanish for to grow old, then say \emph{la vejez} again
\end{itemize}

\begin{tcolorbox}[breakable,colback=teal!4,colframe=teal!35!black,title={Before you move on}]
What does la vejez mean? (\textquotedblleft{}Old age\textquotedblright{}.) Say it once more.
\end{tcolorbox}
\section[anciano]{anciano — elderly; the elderly person}

\label{lesson:ES-C555-anciano}



\begin{quote}
Before the new one: say the Spanish for to grow old, then the Spanish for old age.
\end{quote}

\subsection*{You'll want to know: anciano}
\textbf{anciano} — \textquotedblleft{}elderly; the elderly person\textquotedblright{}.

\begin{quote}
Un anciano me ayudó a encontrar la calle.
\end{quote}

\begin{quote}
\emph{An elderly man helped me find the street.}
\end{quote}

More respectful than \emph{viejo}.

\begin{cousinweb}
\emph{anciano} comes from spoken Latin \emph{*anteānus} 'one from before' (EtymDB), built on \emph{ante} 'before'.
\end{cousinweb}

\begin{grammarlens}[title={What you've built}]
One more word for telling, explaining and describing at B1.
\end{grammarlens}

\subsection*{Your turn}
\begin{itemize}
\raggedright
  \item \emph{Say it:} \emph{anciano}
  \item \emph{Say it:} \emph{anciano}, once more
  \item \emph{Say it:} a sentence of your own with \emph{anciano}
  \item \emph{Recall:} say the Spanish for to grow old, then the Spanish for old age, then say \emph{anciano} again
\end{itemize}

\begin{tcolorbox}[breakable,colback=teal!4,colframe=teal!35!black,title={Before you move on}]
What does anciano mean? (\textquotedblleft{}Elderly; the elderly person\textquotedblright{}.) Say it once more.
\end{tcolorbox}
\section[la muerte]{la muerte — death}

\label{lesson:ES-C555-muerte}



\begin{quote}
Before the new one: say the Spanish for old age, then the Spanish for elderly.
\end{quote}

\subsection*{You'll want to know: la muerte}
\textbf{la muerte} — \textquotedblleft{}death\textquotedblright{}.

\begin{quote}
La muerte de su madre le cambió la vida.
\end{quote}

\begin{quote}
\emph{His mother's death changed his life.}
\end{quote}

\begin{cousinweb}
\emph{muerte} is inherited from Latin \emph{mortem}, the object form of \emph{mors} 'death' (EtymDB). The stressed \emph{o} became \emph{ue}, as in \emph{muerto}.
\end{cousinweb}

\begin{grammarlens}[title={What you've built}]
One more word for telling, explaining and describing at B1.
\end{grammarlens}

\subsection*{Your turn}
\begin{itemize}
\raggedright
  \item \emph{Say it:} \emph{la muerte}
  \item \emph{Say it:} \emph{la muerte}, once more
  \item \emph{Say it:} a sentence of your own with \emph{la muerte}
  \item \emph{Recall:} say the Spanish for old age, then the Spanish for elderly, then say \emph{la muerte} again
\end{itemize}

\begin{tcolorbox}[breakable,colback=teal!4,colframe=teal!35!black,title={Before you move on}]
What does la muerte mean? (\textquotedblleft{}Death\textquotedblright{}.) Say it once more.
\end{tcolorbox}
\section[muerto]{muerto — dead; the dead person}

\label{lesson:ES-C555-muerto}



\begin{quote}
Before the new one: say the Spanish for elderly, then the Spanish for death.
\end{quote}

\subsection*{You'll want to know: muerto}
\textbf{muerto} — \textquotedblleft{}dead; the dead person\textquotedblright{}.

\begin{quote}
El accidente dejó dos muertos.
\end{quote}

\begin{quote}
\emph{The accident left two people dead.}
\end{quote}

Used with \emph{estar}: \emph{Está muerto.} \emph{Muerto de hambre} — starving.

\begin{cousinweb}
\emph{muerto} is inherited from Latin \emph{mortuus} 'dead' (EtymDB), the participle of \emph{morī} 'to die', through spoken Latin \emph{*mortus}.
\end{cousinweb}

\begin{grammarlens}[title={What you've built}]
One more word for telling, explaining and describing at B1.
\end{grammarlens}

\subsection*{Your turn}
\begin{itemize}
\raggedright
  \item \emph{Say it:} \emph{muerto}
  \item \emph{Say it:} \emph{muerto}, once more
  \item \emph{Say it:} a sentence of your own with \emph{muerto}
  \item \emph{Recall:} say the Spanish for elderly, then the Spanish for death, then say \emph{muerto} again
\end{itemize}

\begin{tcolorbox}[breakable,colback=teal!4,colframe=teal!35!black,title={Before you move on}]
What does muerto mean? (\textquotedblleft{}Dead; the dead person\textquotedblright{}.) Say it once more.
\end{tcolorbox}
\section[el difunto]{el difunto — the deceased, the late}

\label{lesson:ES-C555-difunto}



\begin{quote}
Before the new one: say the Spanish for death, then the Spanish for dead.
\end{quote}

\subsection*{You'll want to know: el difunto}
\textbf{el difunto} — \textquotedblleft{}the deceased, the late\textquotedblright{}.

\begin{quote}
La familia del difunto llegó a la iglesia.
\end{quote}

\begin{quote}
\emph{The family of the deceased arrived at the church.}
\end{quote}

Formal. As an adjective: \emph{mi difunto marido} — my late husband.

\begin{cousinweb}
\emph{difunto} comes from Latin \emph{dēfūnctus}, the participle of \emph{dēfungī} 'to finish, to discharge' (Lewis \& Short). The dead are those who have finished their life's work.
\end{cousinweb}

\begin{grammarlens}[title={What you've built}]
One more word for telling, explaining and describing at B1.
\end{grammarlens}

\subsection*{Your turn}
\begin{itemize}
\raggedright
  \item \emph{Say it:} \emph{el difunto}
  \item \emph{Say it:} \emph{el difunto}, once more
  \item \emph{Say it:} a sentence of your own with \emph{el difunto}
  \item \emph{Recall:} say the Spanish for death, then the Spanish for dead, then say \emph{el difunto} again
\end{itemize}

\begin{tcolorbox}[breakable,colback=teal!4,colframe=teal!35!black,title={Before you move on}]
What does el difunto mean? (\textquotedblleft{}The deceased, the late\textquotedblright{}.) Say it once more.
\end{tcolorbox}
